\homework{6}{Mon 10 Oct 2022 20:22}{}
Herstein, section 3.3, page 123: Questions 2, 3, 7
\begin{enumerate}
	\item 2. If $\sigma$ is a $k$-cycle, show that $\sigma$ is an odd permutation if $k$ is even, and is an even permutation if $k$ is odd.
\begin{proof}
	We verify that \[
		(1,2,\ldots,k) = (1,2)(2,3)\ldots(k-1,k)
	.\] 
	Thus a $k$-cycle can be written as product of $k-1$ transpositions. $k$ and $k-1 $ always have opposite parity, so the claim is proved.
\end{proof}

	\item 3. Prove that $\sigma$ and $\tau^{-1} \sigma \tau$, for any $\sigma, \tau \in S_n$, are of the same parity.
\begin{proof}
	Suppose the parity of $\sigma$ is $m$ and the parity of $\tau$ is $n$. Then the parity of $\tau^{-1} \sigma\tau$ is $n + m + n \equiv 2 n + m \equiv m \pmod 2$.
\end{proof}

	\item 7. Show that every element in $A_n$ is a product of $n$-cycles.
\begin{proof}
	Using the identity
	\[
		(1,2,\ldots, k)(k, k-1, \ldots, 3, 1,2) = (1,3,2)
	\] 
	, we can write any $3$-cycle as a product of two $k$-cycles where $k$ is larger or equal to $3$. 

	Since $A_n$ is generated by $3$-cycles, we can write each of the $3$-cycle as two $n$-cycles, and obtain the desired representation.

	 We have just proved the result for $n \ge 3$. For $n = 1$ or $n=2$, $A_n$ become trivial, containing only the identity element which can be represented as $(1,2,\ldots,n)(n,n-1,\ldots,1)$ for any $n \ge 1$.
\end{proof}

\end{enumerate}
Herstein, section 6.1, page 221: Questions 1, 2, 8
\begin{enumerate}[resume]
	\item 1. Prove that if $n>2$, the center of $S_n$ is $(e)$.
\begin{proof}
	Suppose for the sake of contradiction that there exists  $\sigma \in Z(S_n)$ such that $\sigma\neq e$. Let \[
		A := \{x \in [n]: \sigma(x) \neq  x\} 
	.\] 
	Since $\sigma\neq e$, $|A| \ge 2$. Take $a, b \in A$ where $a \neq b$. We claim that $\sigma(a,b) \neq (a,b)\sigma$. 

	Suppose $\sigma(a,b) = (a,b)\sigma$, then $(a,b)^{-1} \sigma (a,b) = \sigma$. Thus if we interchange $a$ and $b$ in the cycle notation of $\sigma$, $\sigma$ remains unchanged. This can happen only when $\sigma = (a,b)\tau$ where $\tau$ is disjoint from $(a,b)$. Suppose $\tau$ is nontrivial, and let $c$ be an element changed by $\tau$. Now $\sigma(a,c) = (a,c) \sigma$, so
	\begin{align*}
		&(a,b)\tau(a,c) = (a,c)(a,b)\tau\\
		&\iff (a,b)(a,c)(a,b)\tau(a,c) = \tau \\
		&\iff (a,b)(a,c)\tau(a,b) (a,c) = \tau \\
		&\iff (a,c,b)\tau(a,c,b) = \tau
	.\end{align*}
	Now $(a,c,b)\tau(a,c,b)(b) =  (a,c,b)\tau(a) = (a,c,b)(a) = c$, but this is a contradiction since this implies $\tau(b) = c$, but $\tau(b) = b$. Therefore $\tau$ is trivial, and $\sigma = (a,b)$.

	However, we have $(a,b)(a,c) \neq (a,c)(a,b)$, so $\sigma \not\in Z(S_n)$, a contradiction. Hence in all cases $\sigma \not\in Z(S_n)$, so $Z(S_n)$ having an element other than $e$ leads to contradiction, so $Z(S_n) = (e)$.
\end{proof}

	\item 2. Prove that if $n>3$, the center of $A_n$ is (e).
\begin{proof}
	Suppose for contradiction that center of $A_n$ is not $(e)$. Take nontrivial $\sigma \in Z(A_n)$. Take  $a,b$ such that $\sigma(a)\neq a$ and $\sigma(b)\neq b$. Because $\sigma(a,b) = (a,b)\sigma$, we have  $(a,b)^{-1}\sigma(a,b) = \sigma$, so $\sigma$ must be in the form $(a,b)\tau$ where $\tau$ is disjoint from $(a,b)$. Since $\sigma \in A_n$, $\sigma$ must be of the form $(a,b)(c,d)\tau$ where $\tau \in A_n$. Now
	\begin{align*}
		&(a,b)(c,d)\tau(a,b,c) = (a,b,c)(a,b)(c,d)\tau\\
		&\iff (a,b,d)\tau(a,b,c) = \tau \\
	.\end{align*}
	But now $(a,b,d)\tau(a,b,c)(a) = (a,b,d)(b) = d$, which means $\tau(a)\neq a$, a contradiction since $\tau$ is disjoint from $a $.

	We have arrived at a contradiction, so the center of $A_n$ must be $(e)$.
\end{proof}

	\item 8. If $M \triangleleft N$ and $N \triangleleft G$, show that $a M a^{-1}$ is normal in $N$ for every $a \in G$.
\begin{proof}
	Take arbitary $n \in N$. Now $a^{-1} n^{-1} a M a^{-1}  n a = (a^{-1} n a)^{-1} M (a^{-1} n a)$. Since $N \triangleleft G$, we have $a^{-1} n a \in N$. Since $M \triangleleft N$, $(a^{-1} n a)^{-1} M (a^{-1} n a) = M$.

	Thus $a^{-1} n^{-1} a M a^{-1}  n a  = M$, hence $n^{-1} (a M a^{-1}) n = a M a^{-1} $, therefore $a M a^{-1}  \triangleleft N$.
\end{proof}

\end{enumerate}
